\section{Bamberg-Modul: Ptolemäus--Walter--Morley für Primzahlvierlinge}

Wir betrachten Primzahlvierlinge der Form
\[
Q(p)=(p,p+2,p+6,p+8),
\]
wobei alle vier Zahlen prim sind. Für $p>3$ tragen diese vier Zahlen im hier verwendeten Sprachschema die Klassen
\[
A,\;B,\;C,\;E.
\]
Die leitende Idee ist dabei, dass $E$ als Rahmen- oder Referenzpol fungiert, während $A,B,C$ die triadische Binnenstruktur bilden. Die Viererordnung ist daher keine bloße Vierermenge, sondern eine \emph{$1+3$-Struktur}.

\subsection{EABC-Einbettung}

Wir verwenden die ebene EABC-Einbettung
\[
E \mapsto (1,0),\qquad
A \mapsto (0,1),\qquad
B \mapsto (-1,0),\qquad
C \mapsto (0,-1),
\]
mit logarithmischen Radien
\[
A_p=(0,\log p),\qquad B_p=(-\log(p+2),0),
\]
\[
C_p=(0,-\log(p+6)),\qquad E_p=(\log(p+8),0).
\]
Damit erhält jeder Primzahlvierling eine klassengeometrische Gestalt.

\subsection{Erste Ebene: ptolemäischer Strukturtyp}

\begin{definition}[Ptolemäischer Defekt]
Für vier Punkte $X_1,X_2,X_3,X_4\in\mathbb R^2$ definieren wir
\[
\Pi(X_1,X_2,X_3,X_4)
:= d_{13}d_{24}-d_{12}d_{34}-d_{23}d_{41},
\qquad d_{ij}=\|X_i-X_j\|.
\]
Der Wert $\Pi$ misst die Abweichung von einer ptolemäisch balancierten Viererfigur.
\end{definition}

Wir unterscheiden zwei innere Anordnungen desselben Vierlings:
\[
Q_{\mathrm I}=(A_p,B_p,C_p,E_p),
\qquad
Q_{\mathrm{II}}=(A_p,C_p,B_p,E_p).
\]
Dazu gehören die Defekte
\[
\Pi_{\mathrm I}(p):=\Pi(A_p,B_p,C_p,E_p),
\qquad
\Pi_{\mathrm{II}}(p):=\Pi(A_p,C_p,B_p,E_p).
\]

\begin{proposition}[Numerischer Hauptbefund]
Für die getesteten Primzahlvierlinge bis $4700$ zeigt sich numerisch:
\[
\Pi_{\mathrm I}(p)\to 0,
\qquad
\Pi_{\mathrm{II}}(p)<0,
\qquad
|\Pi_{\mathrm{II}}(p)|\sim 4(\log p)^2.
\]
Die natürliche Form $Q_{\mathrm I}$ ist also asymptotisch ptolemäisch kompatibel, während die Kreuzform $Q_{\mathrm{II}}$ einen systematischen Strukturdefekt trägt.
\end{proposition}

Dies motiviert die Definition
\[
\Delta E_{EABC}(p):=|\Pi_{\mathrm{II}}(p)|
\]
als \emph{EABC-Strukturspannung}. Sie misst die Nichtverträglichkeit von $E$-Pol und $ABC$-Triade.

\subsection{Zweite Ebene: Walter-Besetzungstyp}

Die kompatible Viererordnung $ABCE$ besitzt zwei natürliche Triangulierungen:
\[
T_{ABC}=\triangle ABC,\qquad T_{ACE}=\triangle ACE,
\]
\[
T_{ABE}=\triangle ABE,\qquad T_{BCE}=\triangle BCE.
\]

Nach Marion Walter ist die Fläche des zugehörigen zentralen Hexagons für jedes Dreieck proportional zu dessen Fläche; im hier verwendeten Normalisierungsschema genügt es daher, die Dreiecksflächen selbst zu vergleichen.

\begin{definition}[Walter-Splits]
Wir definieren
\[
\Delta W_{AC}:=W(T_{ACE})-W(T_{ABC}),
\qquad
\Delta W_{BE}:=W(T_{BCE})-W(T_{ABE}),
\]
wobei $W(T)$ proportional zur Fläche von $T$ ist.
\end{definition}

In der EABC-Einbettung ergeben sich explizit
\[
\Delta W_{AC}
= \frac{1}{20}\bigl(\log p+\log(p+6)\bigr)\bigl(\log(p+8)-\log(p+2)\bigr),
\]
\[
\Delta W_{BE}
= \frac{1}{20}\bigl(\log(p+2)+\log(p+8)\bigr)\bigl(\log(p+6)-\log p\bigr).
\]

\begin{proposition}[Walter-Befund]
Für alle getesteten Primzahlvierlinge gilt numerisch
\[
\Delta W_{AC}>0,
\qquad
\Delta W_{BE}>0,
\]
und beide Größen gehen für $p\to\infty$ gegen $0$.
\end{proposition}

Walter misst damit die \emph{Besetzungsasymmetrie} der kompatiblen Form: der $ACE$-Kanal ist flächenmäßig stets etwas schwerer als $ABC$, und der $BCE$-Kanal etwas schwerer als $ABE$.

\subsection{Dritte Ebene: Morley-Gestalttyp}

Für jedes der vier kanonischen Dreiecke betrachten wir das Morley-Dreieck.

\begin{definition}[Morley-Quotient]
Für ein Dreieck $T$ sei
\[
M(T):=\frac{\operatorname{Area}(\text{Morley-Dreieck von }T)}{\operatorname{Area}(T)}.
\]
Dies ist ein dimensionsloser Gestaltparameter.
\end{definition}

Numerisch zeigt sich für die vier kanonischen Dreiecke
\[
T\in\{T_{ABC},T_{ACE},T_{ABE},T_{BCE}\}
\]
die Konvergenz
\[
M(T)\longrightarrow M_\infty,
\qquad
M_\infty\approx 0.0310889132455,
\]
also auf den Morley-Wert des gleichschenklig-rechtwinkligen Referenzdreiecks.

\begin{definition}[Morley-Defekte und Morley-Splits]
Wir setzen
\[
\delta_M(T):=M(T)-M_\infty,
\]
und definieren
\[
\Delta M_{AC}:=M(T_{ACE})-M(T_{ABC}),
\qquad
\Delta M_{BE}:=M(T_{BCE})-M(T_{ABE}).
\]
\end{definition}

\begin{proposition}[Morley-Befund]
Für alle getesteten Primzahlvierlinge gilt numerisch
\[
\Delta M_{AC}>0,
\qquad
\Delta M_{BE}>0,
\]
und beide Größen gehen für $p\to\infty$ gegen $0$.
\end{proposition}

Morley misst damit die \emph{Gestaltreife} der kanonischen Dreiecke. Während Walter vor allem die Verteilung der Flächenmasse erfasst, misst Morley die Annäherung an eine asymptotische Standardgestalt.

\subsection{Gesamttaxonomie}

Damit ergibt sich für jeden Primzahlvierling $Q(p)$ der Klassifikationsvektor
\[
\mathcal K(Q(p))
=
\bigl(
\Delta E_{EABC}(p),
\Delta W_{AC}(p),
\Delta W_{BE}(p),
\Delta M_{AC}(p),
\Delta M_{BE}(p)
\bigr).
\]

Die fünf Komponenten bedeuten:
\begin{align*}
\Delta E_{EABC}&\quad\text{: ptolemäische Strukturspannung},\\
\Delta W_{AC},\Delta W_{BE}&\quad\text{: Walter-Besetzungsasymmetrien},\\
\Delta M_{AC},\Delta M_{BE}&\quad\text{: Morley-Gestaltasymmetrien}.
\end{align*}

\subsection{Bamberg-Axiom}

\begin{quote}
\textbf{B-WM-1.} Stabile Viererordnungen minimieren die EABC-Strukturspannung und treiben zugleich ihre kanonischen Triangulierungen in eine asymptotische Morley-Standardgestalt.
\end{quote}

Formal:
\[
\Delta E_{EABC}\to \min,
\qquad
\delta_M(T)\to 0
\]
für die kanonischen Teil-Dreiecke $T$.

\subsection{Beispieltabelle}

\begin{center}
\begin{tabular}{r|r r r r r}
$p$ & $\Delta E_{EABC}$ & $\Delta W_{AC}$ & $\Delta W_{BE}$ & $\Delta M_{AC}$ & $\Delta M_{BE}$\\
\hline
$5$    & $18.560817$  & $0.037335$ & $0.050064$ & $0.002414$ & $0.004716$\\
$11$   & $28.985179$  & $0.028491$ & $0.034665$ & $0.001332$ & $0.001843$\\
$101$  & $86.623397$  & $0.005433$ & $0.005645$ & $0.000125$ & $0.000129$\\
$821$  & $180.385148$ & $0.000850$ & $0.000855$ & $0.000011$ & $0.000011$\\
$3461$ & $265.720484$ & $0.000291$ & $0.000292$ & $0.000002$ & $0.000002$\\
\end{tabular}
\end{center}

Die Tabelle zeigt den typischen Trend: Die globale EABC-Strukturspannung der Kreuzform bleibt groß und wächst im gewählten Modell skalenhaft, während die Walter- und Morley-Splits der kompatiblen Form klein werden und gegen $0$ tendieren.
